\documentclass{article}
\usepackage{fuzz}
\begin{document}

\section{Basic Types}

\begin{zed}
  [NAME]
\end{zed}

\section{State}

\begin{schema}{Counter}
  count : \nat
\where
  count \leq 100
\end{schema}

\section{Operations}

An \emph{operation} describes how the system state changes.
The $\Delta$ convention includes both the before-state
and the after-state (primed variables).

\begin{schema}{Increment}
  \Delta Counter
\where
  count < 100 \\
  count' = count + 1
\end{schema}

The precondition $count < 100$ ensures the invariant is
preserved after incrementing.

\begin{schema}{Reset}
  \Delta Counter
\where
  count' = 0
\end{schema}

\texttt{Reset} has no precondition --- it succeeds in any state.

A $\Xi$ operation observes the state without changing it.

\begin{schema}{GetCount}
  \Xi Counter \\
  result! : \nat
\where
  result! = count
\end{schema}

The \texttt{!} suffix marks an output variable.

\end{document}
